% =====================================================================
% SECTION 1: Introduction and method  +  SECTION 2: Arithmetic sector
% =====================================================================

\section*{Abstract}
\noindent
This document is a rigorous reconstruction of the ``prime-spectral framework'' developed in a
forty-eight-turn source conversation, in which the statistical mechanics of a free Bose gas of
``primons'' (one-particle energies $\log p$, $p$ ranging over the rational primes) was advanced as
the arithmetic core of a unified picture of quantum gravity, time, and the arrows of time. Our
procedure is threefold. First, we \emph{elevate} every statement that admits a precise formulation:
the primon Hilbert space is defined as $\ell^{2}$ of the free commutative monoid on the primes,
sidestepping the nonseparability of naive infinite tensor products; the primon Hamiltonian is
proved self-adjoint with simple spectrum $\hbar\omega_{0}\log n$; the partition function is
identified with the Riemann zeta function; the exponential state-counting law
$N(E)=\lfloor e^{E/\hbar\omega_{0}}\rfloor$ is given its correct name, a \emph{Hagedorn spectrum}
with critical temperature $T_{H}=\hbar\omega_{0}/k_{B}$; and the ``thermodynamic cascade into
composite states'' is grounded in the Hardy--Ramanujan and Erd\H{o}s--Kac theorems. Second, we
\emph{repair} the central gaps: the frozen-Hamiltonian problem (the diagonal primon Hamiltonian
admits no transitions), the false identification of composite integers with entangled states
(occupation-basis states are product states across every prime-mode bipartition), and the
overstrong no-recurrence claim (exact periodicity is indeed forbidden by the
$\Qq$-linear independence of prime logarithms, but approximate recurrence survives in every
finite truncation). Third, we \emph{demote} to an explicit, falsifiable conjecture register every
claim whose proof is out of reach, from the zeta/black-hole partition-function identification to
the Galois--Langlands dictionary for the Standard Model. The result is a framework whose
mathematical core is now theorem-grade, whose analogies are tightened to citable results, and
whose remaining ambitions are stated with enough precision to be attacked or refuted.

\bigskip

\section{Introduction: Source, Method, and Results}\label{sec:intro}

\subsection{The source conversation}

The source material is a shared conversation (48 turns) in which a user and an assistant
co-developed a speculative but internally ambitious framework: that the prime numbers, through the
primon gas of Julia~\cite{julia1990} and Spector~\cite{spector1998}, supply the quantum-mechanical
spectrum of the near-singularity gravitational field, that the Riemann zeta function is the
gravitational partition function, and that the entire architecture of time---its problem, its
arrow, its block-universe metaphysics, and its impossibility of cycles---follows from the
arithmetic of the integers. The conversation drew on reported 2025--2026 developments linking
conformal primon gases to black-hole interiors~\cite{hartnollyang2026}, on the philosophical
literature of time's arrow~\cite{carroll2010,wallace2017,price1996}, and on first-person
metaphysics~\cite{hellie2013,hare2009,conitzer2019}. It culminated in claims that Gaussian primes
extend the framework to five dimensions, that the Absolute Galois group plays the role of a grand
unified symmetry, and that gravity is an entropic, $1/N$-suppressed collective effect.

Taken verbatim, the conversation alternates between genuine mathematics, physics speculation, and
rhetorical overreach. That alternation is precisely what makes it worth reconstructing: the
skeleton is sound and unusually coherent, while the joints are loose. The present document
keeps the skeleton, tightens every joint, discards nothing that can be saved by honest
reformulation, and labels everything else.

\subsection{Method: elevate, repair, demote}\label{sec:method}

Every substantive claim in the source was subjected to one of three treatments.

\begin{enumerate}[leftmargin=1.6em]
\item \textbf{Elevate.} If a claim has a precise mathematical formulation and a proof, it is
stated as a theorem, lemma, or proposition with its proof or an exact citation. Examples: the
$\Qq$-linear independence of $\{\log p\}$; the identification of the partition function with
$\zeta$; the Hagedorn asymptotics; the positive-definiteness of the Page--Wootters clock kernel;
the exact state-counting identity $\log N(E)=E/\hbar\omega_{0}+O(1)$.
\item \textbf{Repair.} If a claim is right in spirit but wrong in mechanism, the mechanism is
replaced and the repaired claim is re-proved. Three repairs dominate this document:
(i)~the diagonal primon Hamiltonian cannot drive the claimed ``cascade into composite states'',
so a multiplicative transition sector (multiplication operators $M_{n}$ and a hopping dynamics
on the \emph{prime lattice}) is introduced in Section~\ref{sec:dynamics};
(ii)~composite integers are \emph{not} entangled states---the occupation basis is a product
basis---so arithmetic entanglement is redefined across prime-mode bipartitions in
Section~\ref{sec:entangle}, and the ``arithmetic holism'' that resolves Bell's trilemma is
relocated from single integers to superpositions;
(iii)~the ``universe cannot repeat'' claim conflated exact periodicity, Poincar\'e recurrence,
and dynamical return; these are separated in Section~\ref{sec:recurrence} into one theorem and
two counterexamples.
\item \textbf{Demote.} If a claim is plausible, well-formed, and beyond current proof
technology, it is restated as a numbered conjecture with evidence, dependencies, and a
falsifier. Twelve conjectures ($\mathsf{C1}$--$\mathsf{C12}$) are consolidated in
Section~\ref{sec:register}. The guiding rule is that a conjecture must be \emph{attackable}:
each one carries at least one specified route to confirmation or refutation.
\end{enumerate}

An audit trail of the corrections is given in Appendix~\ref{app:corrections}; it lists each
original claim, its repaired or demoted status, and the pointer to the section where the
treatment occurs. Nothing in the source ambition is silently dropped: where a claim is rejected
outright (for example, that the Hamiltonian $H_{P}=\hbar\omega_{0}\log\Nhat$ by itself ``drives
transitions into composite integer states''), the rejection is argued, and the surviving
content is carried forward under its own name.

\subsection{Summary of results}

The reconstruction yields a layered structure. The \emph{mathematical floor} (all theorem-grade)
consists of: the arithmetic Hilbert space $\Hil_{\mathcal{A}}=\ell^{2}(\Nn,\times)$
(Section~\ref{sec:arithmetic}); the self-adjoint primon Hamiltonian with simple spectrum
$\hbar\omega_{0}\log n$; the partition function $Z(\beta)=\zeta(\beta\hbar\omega_{0})$ with
Hagedorn temperature $T_{H}=\hbar\omega_{0}/k_{B}$; the exact counting law
$\log N(E)=E/\hbar\omega_{0}+O(1)$; the Erd\H{o}s--Kac statistics of the microcanonical
ensemble; the isometric representation $n\mapsto M_{n}$ of the multiplicative monoid on
$\Hil_{\mathcal{A}}$; the phase-incommensurability theorem forbidding exact periodicity; the
Bohr almost-periodicity of the zeta clock's correlation kernel; and the product-state theorem
that repairs the entanglement story. The \emph{conjectural superstructure} consists of the
twelve register entries, of which the deepest are $\mathsf{C1}$ (the zeta/black-hole
identification), $\mathsf{C3}$ (the arithmetic holographic bound), $\mathsf{C4}$ (thermalization
on the prime lattice), and $\mathsf{C9}$ (the Galois--gauge dictionary).

The philosophical theses of the source---eternalism with perspectival observers, an extrinsic
arrow grounded in arithmetic state-counting, derivative causality, and the rejection of
superdeterminism---survive the reconstruction in \emph{conditional} form: they are
well-defined consequences of the conjectural superstructure, and they are marked as
interpretive rather than theorem-grade wherever the mathematics cannot carry them
(Section~\ref{sec:philosophy}).

\section{The Arithmetic Sector}\label{sec:arithmetic}

This section fixes the foundational objects. Everything here is elementary but must be stated
carefully, because the source conversation (and part of the physics literature) informally
writes the primon Fock space as an infinite tensor product, an object that is \emph{not}
separable and is the wrong home for the theory. The repair is cheap and clarifying.

\subsection{The arithmetic Hilbert space}

Let $\Pr=\{2,3,5,\dots\}$ denote the rational primes. Let
\begin{equation}\label{eq:monoid}
\mathcal{F}\;=\;\Bigl\{k\colon \Pr\to\Nn \;:\; k_{p}\neq 0 \text{ for only finitely many } p\Bigr\}
\end{equation}
be the free commutative monoid on the primes. By the fundamental theorem of arithmetic, the
map $\mathcal{F}\to\Nn_{\geq 1}$, $k\mapsto \prod_{p}p^{k_{p}}$, is a monoid isomorphism
between $(\mathcal{F},+)$ and $(\Nn_{\geq 1},\times)$, where $+$ denotes componentwise addition
of occupation vectors. This single classical theorem is the load-bearing wall of the entire
framework: it is what makes ``states of the arithmetic universe'' and ``positive integers''
the same bookkeeping.

\begin{definition}[Arithmetic Hilbert space]\label{def:hilbert}
The \emph{arithmetic Hilbert space} is
\begin{equation}
\Hil_{\mathcal{A}}\;=\;\ell^{2}(\mathcal{F})\;\cong\;\ell^{2}(\Nn_{\geq 1}),
\end{equation}
with orthonormal basis $\{\ket{k}\}_{k\in\mathcal{F}}$, written $\ket{n}$ for
$n=\prod_{p}p^{k_{p}}$ when no confusion arises. Each prime $p$ defines a \emph{mode} with
occupation counter
\begin{equation}
\hat n_{p}\ket{k}=k_{p}\ket{k},
\end{equation}
a self-adjoint operator with spectrum $\Nn$.
\end{definition}

\begin{remark}[Why not an infinite tensor product?]
The source writes $\Fock_{\Pr}=\bigotimes_{p\in\Pr}\Fock_{p}$ with one bosonic oscillator per
prime. The literal von Neumann infinite tensor product of the $\ell^{2}(\Nn)$ over a countably
infinite index set is nonseparable, and its physically relevant separable subspaces depend on
a choice of reference vector. Definition~\ref{def:hilbert} achieves the intended object
directly: $\ell^{2}$ of the countable basis $\mathcal{F}$ is separable, and the tensor
factorization over any \emph{finite} set of modes is available whenever needed
(Proposition~\ref{prop:productstructure}). This is the first, and easiest, instance of the
general principle of this reconstruction: replace informal infinite constructions by the
countable objects they were meant to describe.
\end{remark}

For each prime $p$ the standard bosonic ladder operators act on the $p$-mode: on basis vectors
\begin{equation}
a_{p}^{\dagger}\ket{k}=\sqrt{k_{p}+1}\,\ket{k+e_{p}},
\qquad
a_{p}\ket{k}=\sqrt{k_{p}}\,\ket{k-e_{p}},
\end{equation}
where $e_{p}$ is the unit occupation vector, and $[a_{p},a_{q}^{\dagger}]=\delta_{pq}$ holds
on the dense span of basis vectors. The total particle number is
$\Nhat_{\mathrm{tot}}=\sum_{p}\hat n_{p}$, self-adjoint on
$\{\psi:\sum_{p}k_{p}^{2}|\psi_{k}|^{2}<\infty\}$.

\subsection{The multiplicative number operator and the primon Hamiltonian}

The source's central operator is the \emph{multiplicative number operator}, whose eigenvalue on
$\ket{k}$ is the integer $n=\prod_{p}p^{k_{p}}$ itself. We define it together with the
Hamiltonian it generates.

\begin{definition}[Multiplicative number operator; primon Hamiltonian]\label{def:hamiltonian}
The \emph{multiplicative number operator} $\Nhat$ is the diagonal operator
\begin{equation}
\Nhat\ket{k}=\Bigl(\prod_{p}p^{k_{p}}\Bigr)\ket{k},
\end{equation}
self-adjoint on its maximal domain. The \emph{primon Hamiltonian} is
\begin{equation}\label{eq:primonH}
H_{P}\;=\;\hbar\omega_{0}\log\Nhat
\;=\;\hbar\omega_{0}\sum_{p\in\Pr}(\log p)\,\hat n_{p},
\qquad \omega_{0}>0,
\end{equation}
where the equality of the two expressions is pointwise on basis vectors and extends by
closure.
\end{definition}

Equation~\eqref{eq:primonH} exhibits $H_{P}$ as a \emph{free Bose gas} whose one-particle
spectrum is $\{\log p:p\in\Pr\}$: each prime is one species of elementary excitation (a
\emph{primon}), and a many-body state $\ket{n}$ has energy equal to the logarithm of the
integer it represents. This is Julia's primon gas~\cite{julia1990}, cleaned of its informal
packaging.

\begin{theorem}[Self-adjointness, spectrum, and simplicity]\label{thm:spectrum}
$H_{P}$ is self-adjoint and nonnegative. Its spectrum is the pure point set
\begin{equation}
\spec(H_{P})=\overline{\{\hbar\omega_{0}\log n:n\in\Nn_{\geq 1}\}},
\end{equation}
a discrete set whose only accumulation point is $+\infty$; it has no finite accumulation
point. Every eigenvalue $\hbar\omega_{0}\log n$ is
\emph{simple}: $\log m=\log n$ if and only if $m=n$. The ground state is the vacuum
$\ket{1}$ with eigenvalue $0$, and it is unique.
\end{theorem}

\begin{proof}
Under the unitary $U\colon\Hil_{\mathcal{A}}\to\ell^{2}(\Nn_{\geq 1})$,
$U\ket{k}=\ket{n}$ (unitary by the fundamental theorem of arithmetic),
$H_{P}$ becomes multiplication by the real sequence
$a_{n}=\hbar\omega_{0}\log n$. Multiplication by a real-valued sequence is self-adjoint on its
maximal domain $\{c\in\ell^{2}:\sum_{n}a_{n}^{2}|c_{n}|^{2}<\infty\}$, and its spectrum is the
closure of $\{a_{n}\}$. Since $a_{n}\to\infty$ monotonically, the closure adds no finite
points. Simplicity: if $a_{m}=a_{n}$ with $m\neq n$, then $\log m=\log n$ forces $m=n$, a
contradiction. The ground state: $a_{n}\geq 0$ with equality only at $n=1$.
\end{proof}

Two consequences deserve emphasis because the source conversation uses both constantly.
First, the \emph{unique zero-entropy ground state} $\ket{1}$ is the multiplicative identity;
this is the rigorous content of the source's ``arithmetic floor'' and of its Janus point
claims, developed in Sections~\ref{sec:thermo} and~\ref{sec:philosophy}. Second, the
\emph{simplicity} of every level means the primon spectrum has no degeneracy whatsoever: the
integers label states one-to-one. Any degeneracy in a putative physical realization must
therefore come from sectors \emph{beyond} $H_{P}$ (for instance the Gaussian sector of
Section~\ref{sec:galois}).

\subsection{The partition function and the Hagedorn temperature}

\begin{proposition}[Partition function]\label{prop:partition}
For inverse temperature $\beta>0$,
\begin{equation}\label{eq:zetaZ}
Z(\beta)\;=\;\Tr\,e^{-\beta H_{P}}
\;=\;\sum_{n=1}^{\infty}e^{-\beta\hbar\omega_{0}\log n}
\;=\;\sum_{n=1}^{\infty}n^{-\beta\hbar\omega_{0}}
\;=\;\zeta(\beta\hbar\omega_{0}),
\end{equation}
and the series converges if and only if $\beta\hbar\omega_{0}>1$, i.e.
$T<T_{H}$, where
\begin{equation}\label{eq:hagedornT}
T_{H}\;=\;\frac{\hbar\omega_{0}}{k_{B}}
\end{equation}
is the \emph{Hagedorn temperature} of the primon gas. On the half-line $\sigma=\beta\hbar\omega_{0}>1$,
the Euler product
\begin{equation}\label{eq:eulerproduct}
\zeta(\sigma)=\prod_{p\in\Pr}\frac{1}{1-p^{-\sigma}}
=\prod_{p\in\Pr}\frac{1}{1-e^{-\beta\varepsilon_{p}}},
\qquad \varepsilon_{p}:=\hbar\omega_{0}\log p,
\end{equation}
is precisely the grand-canonical factorization of a free Bose gas into independent prime modes,
with each mode contributing the standard Bose factor. The mean occupation of mode $p$ is
\begin{equation}
\langle \hat n_{p}\rangle_{\beta}=\frac{1}{e^{\beta\varepsilon_{p}}-1}
=\frac{1}{p^{\beta\hbar\omega_{0}}-1}.
\end{equation}
\end{proposition}

\begin{proof}
The trace formula is the spectral decomposition of Theorem~\ref{thm:spectrum}; convergence at
$\sigma>1$ and the Euler product are Euler's theorems; the occupation formula is the standard
free-boson computation on each mode.
\end{proof}

The source conversation states the identification $Z=\zeta$ correctly but never names the
consequence that gives it physical teeth: the primon gas is a \emph{Hagedorn system}
\cite{hagedorn1965}. Its density of states grows exponentially in the energy, and the
partition function diverges at a finite critical temperature. In string theory this is the
famous signal that, near $T_{H}$, energy input goes into multiplying the number of light
degrees of freedom rather than into heating. The same reading applies here with total
exactness, because the counting is exactly solvable.

\begin{theorem}[Hagedorn asymptotics]\label{thm:hagedorn}
Let $k_{B}=1$ and write $T_{H}=\hbar\omega_{0}$. As $T\uparrow T_{H}$ (i.e.
$\sigma=\beta\hbar\omega_{0}\downarrow 1$):
\begin{align}
\zeta(\sigma)&=\frac{1}{\sigma-1}+\gamma+O(\sigma-1),\\
U(T):=-\frac{\partial}{\partial\beta}\log Z(\beta)
&=\frac{\hbar\omega_{0}}{\sigma-1}+O(\hbar\omega_{0})
\;=\;\frac{T\,T_{H}}{T_{H}-T}+O(T_{H}),\\
S(T):=\beta U+\log Z
&=\frac{1}{\sigma-1}-\log(\sigma-1)+O(1),\\
C(T):=\frac{dU}{dT}
&=\frac{T_{H}^{2}}{(T_{H}-T)^{2}}\,\bigl(1+O((T_{H}-T)/T)\bigr).
\end{align}
In particular the specific heat diverges as $T\uparrow T_{H}$, while the entropy diverges
only as $T/(T_{H}-T)$; energy is absorbed at nearly constant temperature.
\end{theorem}

\begin{proof}
The Laurent expansion $\zeta(\sigma)=1/(\sigma-1)+\gamma+O(\sigma-1)$ is classical
\cite{titchmarsh1986}. Then
$-\zeta'/\zeta(\sigma)=[1/(\sigma-1)^{2}+O(1)]/[1/(\sigma-1)+O(1)]=1/(\sigma-1)+O(1)$, giving
$U=\hbar\omega_{0}(-\zeta'/\zeta)(\sigma)$. Substituting $\sigma=T_{H}/T$ gives
$\sigma-1=(T_{H}-T)/T$ and the stated form. The entropy follows from
$S=\beta U+\log Z$ with $\log Z=-\log(\sigma-1)+O(1)$. Differentiating $U=T T_{H}/(T_{H}-T)$
in $T$ gives $C=T_{H}^{2}/(T_{H}-T)^{2}$ up to the stated relative error.
\end{proof}

\begin{remark}[Microcanonical reading]
Section~\ref{sec:thermo} proves the exact counting law
$\log N(E)=E/\hbar\omega_{0}+O(1)$, which says that the \emph{microcanonical} temperature
$(dS/dE)^{-1}$ equals $T_{H}$ identically, at every energy scale: the primon gas is an exactly
critical system. It never heats; it only multiplies states. This is the sharpened form of the
source's statement that ``the state space expands exponentially'': the expansion rate is not
merely large, it is \emph{uniform} in $E$, and that uniformity is what the pole of $\zeta$ at
$s=1$ encodes. The classical smooth term ``$x$'' of the prime number theorem is the same
statement read on the counting function $\psi(x)$, and the oscillatory remainder
$-\sum_{\rho}x^{\rho}/\rho$ is the fluctuation sector; the reconstruction returns to this
duality in Sections~\ref{sec:time} and~\ref{sec:galois}.
\end{remark}

\subsection{Gaussian extension}

The source extends the arithmetic sector to the Gaussian integers $\Zz[i]$ for the putative
five-dimensional (and chirality-carrying) sector, and here the mathematics is genuinely
supportive, because the extension is an instance of the standard theory of number fields.

\begin{definition}[Gaussian primon gas]\label{def:gaussian}
Let $\mathcal{P}_{\Zz[i]}$ denote the Gaussian primes modulo units $\{\pm1,\pm i\}$, with norm
$N(\pi)=|\pi|^{2}$. The \emph{Gaussian arithmetic Hilbert space} is
$\Hil_{\Zz[i]}=\ell^{2}(\mathcal{F}_{\Zz[i]})$ over the free commutative monoid on
$\mathcal{P}_{\Zz[i]}$, and the Gaussian primon Hamiltonian is
\begin{equation}
H_{\Zz[i]}=\hbar\omega_{0}\sum_{\pi\in\mathcal{P}_{\Zz[i]}}\log N(\pi)\;\hat n_{\pi},
\qquad
Z_{\Zz[i]}(\beta)=\prod_{\pi}\frac{1}{1-N(\pi)^{-\beta\hbar\omega_{0}}}.
\end{equation}
\end{definition}

\begin{theorem}[Dedekind factorization]\label{thm:dedekind}
$Z_{\Zz[i]}(\beta)=\zeta_{\Qq(i)}(\beta\hbar\omega_{0})$, the Dedekind zeta function of the
Gaussian rationals, and
\begin{equation}\label{eq:factorization}
\zeta_{\Qq(i)}(s)=\zeta(s)\,L(s,\chi_{-4}),
\end{equation}
where $L(s,\chi_{-4})=\sum_{m}(-1)^{m-1}(2m-1)^{-s}$ is the Dirichlet $L$-function of the
nontrivial character mod $4$. Both factors converge for $\sigma>1$; the product is
entire on the common half-plane and $L(1,\chi_{-4})=\pi/4\neq 0$.
\end{theorem}

\begin{proof}
Unique factorization of ideals in $\Zz[i]$ (it has class number $1$) gives the Euler product
over Gaussian prime ideals; identifying prime ideals with Gaussian primes modulo units gives
the stated product, and the splitting classification of rational primes in $\Zz[i]$
($p\equiv1\ (4)$ splits, $p\equiv3\ (4)$ is inert, $2$ ramifies) yields
$\zeta_{\Qq(i)}=\zeta\,L(\chi_{-4})$ \cite{neukirch1999}. The nonvanishing
$L(1,\chi_{-4})=\pi/4$ is Leibniz's series.
\end{proof}

Equation~\eqref{eq:factorization} is the exact arithmetic statement behind the source's slogan
that ``the fifth dimension requires complex primes.'' The Gaussian sector is the rational
sector \emph{times} a chiral factor: $L(s,\chi_{-4})$ is precisely the asymmetry between the
splitting class $p\equiv1\ (4)$ and the inert class $p\equiv3\ (4)$. This is developed into a
precise (and appropriately demoted) arithmetic model of charge conjugation and CP violation in
Section~\ref{sec:galois}.

\subsection{What the arithmetic sector does and does not assert}

It is worth fixing the epistemic boundary before dynamics is introduced. Everything in this
section is unconditional mathematics: the primon gas exists as a quantum system, its spectrum
is exactly the logarithms of the integers, its thermodynamics is exactly the analytic theory
of $\zeta$, and its critical behavior is exactly Hagedorn behavior. \emph{No claim has yet
been made that any physical system is a primon gas.} The identification
$Z^{\mathrm{ren}}_{\mathrm{grav}}(s)=\zeta(s)$ for near-singularity gravity is
Conjecture~$\mathsf{C1}$ of the register, the reported 2025--2026 developments
\cite{hartnollyang2026} being its current physical evidence. The discipline of this
reconstruction is that such identifications live in the register, not in the theorems.
